\documentclass[12pt,a4paper]{article}
\usepackage[utf8]{inputenc}
\usepackage[T1]{fontenc}
\usepackage{lmodern}
\usepackage{graphicx}
\usepackage{setspace}
\usepackage{geometry}

\begin{document}

\begin{center}
    \centering

    {\huge \textbf{Análise, Projeto e Desenvolvimento Ágil}}\\[0.5cm]

    {\Huge \textbf{Questões de Extreme Programming}}\\[0.5cm]

    {\large Aluno 01: \textbf{Henrique Maia Cardosa}}\\[0.5cm]
    {\large Aluno 02: \textbf{João Miguel de Castro Menna}}\\[0.5cm]
    {\large Aluno 03: \textbf{Ricardo Gabriel Fialho Santos}}\\[2.5cm]
\end{center}


\pagebreak

\section*{Questões sobre Extreme Programming (XP)}

\paragraph{}

1. Valores do XP e consequências práticas Explique como Comunicação,
Simplicidade, Feedback, Coragem e Respeito se traduzem em decisões do dia a dia
(ex.: revisão de código, padrão de código, spikes). Dê um exemplo de trade-off onde
“Simplicidade” e “Coragem” evitam over engineering.

R: Os valores do XP orientam o trabalho diário. Comunicação mantém o time alinhado,
simplicidade evita complicações, feedback permite ajustes rápidos, coragem incentiva
mudanças necessárias e respeito sustenta o trabalho em equipe. Simplicidade e coragem
juntos evitam overengineering ao focar apenas no essencial.

\paragraph{}

2. TDD — especificação executável Descreva o ciclo Red → Green → Refactor e
explique por que TDD melhora design e testabilidade. Cite dois anti-padrões comuns
e como evitá-los.

R: O TDD segue o ciclo Red, Green e Refactor: cria o teste que falha, faz o código passar e
melhora o design. Isso gera código limpo e testável. Anti-padrões comuns são testes frágeis
e irrelevantes; evitam-se testando comportamento, não implementação.

\paragraph{}

3. Pair Programming — qualidade e fluxo Quando pair programming traz melhor
custo-benefício? Defina papéis Driver/Navigator, proponha um protocolo de sessão
(tempo, rotatividade, micro-retros) e aponte dois riscos e respectivas mitigações.

R: Pair programming é útil em partes críticas e onboarding. O driver codifica e o navigator
orienta. Sessões curtas com trocas frequentes e pausas evitam fadiga e diferenças de
ritmo.

\paragraph{}

4. CI e releases pequenos — o “cadence técnico” do XP Esboce um pipeline de CI
mínimo (estágios e gates) alinhado ao XP e defina métricas para detectar regressão
de qualidade sem desacelerar o time. Como small releases reduzem risco?

R: O pipeline de CI deve incluir build, testes, lint e deploy. Métricas como tempo de build e
falhas mostram regressões. Releases pequenas reduzem risco e facilitam correções
rápidas.

\paragraph{}

5. Refactoring, DoD técnico e YAGNI (decidir hoje vs. amanhã)
Diferencie refactoring de reescrita; proponha um DoD técnico mínimo para XP (testes,
CI, padrões, observabilidade) e explique quando invocar YAGNI sem comprometer
qualidades como segurança ou operabilidade.

R: Refactoring melhora o design sem mudar o comportamento, diferente da reescrita. O DoD
inclui testes, CI verde e padrões seguidos. O YAGNI evita implementar o que não é preciso
sem comprometer segurança ou estabilidade

\end{document}
